
\section{实验结果与分析}
\label{sec:ch3_experiment}

为验证所提 BC-CTDE-MAPPO 方法在天基计算网络协同资源分配场景中的有效性，本节从收敛性、不同负载水平下的系统性能、规模扩展与参数敏感性、抗自私行为鲁棒性以及信誉机制效果五个方面开展仿真实验。为保证结果的统计可靠性，所有实验均独立运行 5 次，报告均值与标准差。实验指标主要包括任务完成率、平均任务时延、Jain 公平指数、单位任务能耗以及链上信誉演化趋势。为保证对比的完整性，本文选取 MAPPO、MADDPG、IPPO、Greedy 和 Random 作为基线方法。

\subsection{实验设置}
实验场景中包含 20 个异构卫星计算节点，分布在 3 个轨道层上，系统共考虑 4 类在轨任务，训练总回合数为 500，每回合包含 200 个时隙。Actor 与 Critic 的学习率分别设置为 $3\times10^{-4}$ 和 $1\times10^{-3}$，折扣因子取 $\gamma=0.99$，GAE 参数取 $\lambda=0.95$，PPO 裁剪系数取 0.2。混合奖励函数中效率项、公平项与可信项的权重分别设置为 0.5、0.3 和 0.2。联盟链默认同步时延设为 1.0 s，任务到达率从 0.3 增加至 1.1，自私节点比例从 0\% 增加至 40\%。详细参数如表~\ref{tab:exp_setting} 所示。

\subsection{收敛性分析}
图~\ref{fig:3_7}(a) 给出了不同训练算法在 500 回合中的平均回报收敛曲线，阴影区域为 5 次独立运行的 $\pm1\sigma$ 标准差范围，体现了结果的统计可靠性。BC-CTDE-MAPPO 在约 200 回合后已进入稳定收敛区间，最终平均回报（$88.4\pm2.1$）明显高于其余基线方法，且阴影带最窄，表明训练方差最小。图~\ref{fig:3_7}(b) 以误差棒形式进一步比较了不同算法的收敛回合数（误差棒大小根据各算法后 50 回合的回报标准差等比缩放）。BC-CTDE-MAPPO 的收敛回合最少，仅为 210 回合；MAPPO、MADDPG 和 IPPO 分别需要 260、330 和 360 回合。结合表~\ref{tab:ch3_convergence} 可知，BC-CTDE-MAPPO 在最终平均回报、收敛速度及后 50 回合回报标准差等指标上均优于对比方法。

\subsection{不同负载水平下的系统性能分析}
从图~\ref{fig:3_8}(a) 可以看出，随着任务到达率由 0.3 增加到 1.1，所有算法的任务完成率均有所下降，图中误差棒为 5 次独立运行的 $\pm1\sigma$ 标准差。BC-CTDE-MAPPO 在整个负载区间内始终保持最优性能，其任务完成率由 98.2\% 下降至 89.1\%，下降幅度（9.1pp）明显小于 MAPPO（13.5pp）和 Greedy（23.6pp）。值得注意的是，在图中以背景色块标注的高负载区间（$\lambda=0.9\sim1.1$），各方法之间的性能差距显著拉大。低负载时（$\lambda\leq0.5$）所有方法均能维持 90\% 以上的完成率，差距在 4pp 以内；而在 $\lambda=1.1$ 的高负载条件下，BC-CTDE-MAPPO（89.1\%）与 Greedy（70.8\%）相差近 18.3pp，与 Random（58.9\%）相差逾 30pp。这说明本文方法的多智能体协同调度与负载均衡机制的核心价值在高竞争场景中才得到充分体现。图~\ref{fig:3_8}(b) 给出了不同负载水平下的平均任务时延变化趋势，高负载区间同样以背景色块标注，BC-CTDE-MAPPO 的时延增长最缓。表~\ref{tab:ch3_load_perf} 表明 BC-CTDE-MAPPO 在各负载水平下均保持最高 Jain 公平指数。

\subsection{规模扩展与参数敏感性分析}
本文将节点规模从 20 逐步扩展至 40、60 和 80，图~\ref{fig:3_9}(a) 表明 BC-CTDE-MAPPO 的下降最缓（从 92.3\% 降至 86.8\%，降幅 5.5pp）。图~\ref{fig:3_9}(b) 给出了链上附加时延对任务完成率的影响，当延迟位于 0.5--2.0 s 区间时性能下降较为平缓，但超过 2.5 s 后性能显著恶化。

消融实验结果如图~\ref{fig:3_10} 所示，两个子图分别展示了任务完成率和 Jain 公平指数，每个柱子均附有 $\pm1\sigma$ 误差棒（n=5）。图~\ref{fig:3_10}(a) 中，左侧实色柱为全场景均值，右侧斜线柱为 30\% 自私节点场景下的结果。去除区块链公共信息板后任务完成率降至 89.1\%（全场景 $-3.2$pp，30\% 自私节点下 $-6.7$pp），说明区块链作为可信公共信息板的作用在异常节点存在时尤为关键。图~\ref{fig:3_10}(b) 中，去除公平奖励项后 Jain 公平指数从 0.918 大幅降至 0.842（$-0.076$），而去除可信奖励项对公平指数影响相对较小（0.901），但在自私节点场景下完成率降至 84.1\%（$-5.3$pp），证明公平项与可信项分别作用于不同维度，缺一不可。表~\ref{tab:ch3_ablation} 给出了完整数值汇总。

\subsection{抗自私行为鲁棒性分析}
本实验中，自私节点的异常行为模型定义如下：节点以概率 $p_s=0.7$ 拒绝接受来自其他节点的卸载任务（任务拒收），同时以均值 $\mu_f = -0.3$ 的偏移量虚报自身当前负载（虚报负载），在接受任务时优先选择计算量小、收益高的任务（偏好高收益任务）。三类异常行为同时激活，共同模拟机会主义节点在去中心化网络中的典型偏离策略。仿真中通过随机选取一定比例的卫星节点激活上述行为，横轴表示自私节点在总节点中的占比。

随着自私节点比例从 0\% 提高到 40\%，图~\ref{fig:3_9}(a) 中 BC-CTDE-MAPPO 的任务完成率退化最小（降幅 6.0pp），MAPPO 降幅为 11.2pp，所有结果均附有 $\pm1\sigma$ 误差棒（n=5）。图~\ref{fig:3_9}(b) 给出了 Jain 公平指数随自私节点比例变化的曲线，同样附有误差棒。在 30\% 自私节点条件下，BC-CTDE-MAPPO 的 Jain 公平指数仍达到 0.886（$\pm0.011$），而 MAPPO 和 Greedy 分别下降至 0.816 和 0.722，说明链上信誉机制能够有效抑制自私行为对系统公平性的破坏。

\subsection{信誉机制效果分析}
合作节点的信誉值随时间步增加逐渐上升并稳定在 0.90 以上，自私节点信誉值从初始 0.50 持续下降至 0.20 以下，验证了链上信誉反馈机制能够有效区分合作行为与机会主义偏离行为。

\subsection{结果讨论}
综合以上结果，可以得到三点核心结论。其一，可信共享与信誉约束在长期运行中比单步优化更重要。其二，区块链为公共状态共享、行为审计和激励兑现提供了统一执行媒介。其三，本文方法的优势在高负载、强异构和异常节点存在时更加明显。
